\chapter{Fundamentação teórica}
\label{cap:fundamentos}

\section{Linguagens de descrição de hardware e VHDL}

Uma linguagem de descrição de hardware (HDL) é usada para a descrição formal e o projeto de
circuitos eletrônicos, sobretudo digitais. Ela permite descrever o funcionamento, a concepção
e a organização de um circuito e testá-lo por meio de simulação (GARCIA; MAGALHÃES, 2024).
Dentre as vantagens de uma HDL padronizada estão a independência do fornecedor, a porta-
bilidade e a reusabilidade. Em relação ao diagrama esquemático, a descrição por HDL oferece
projeto independente da tecnologia, facilidade de atualização, redução do tempo e do custo
de projeto e redução do volume de documentação (GARCIA; MAGALHÃES, 2024).

\secao{Estrutura de um projeto em VHDL}

Um projeto VHDL possui três unidades fundamentais \cite{aula3}:

\begin{listar}
  \item \texttt{LIBRARY}: lista os pacotes e bibliotecas usados no projeto;
  \item \texttt{ENTITY}: especifica as entradas e saídas do circuito, com o
        nome do arquivo \texttt{.vhd} igual ao nome da entidade;
  \item \texttt{ARCHITECTURE}: contém o código que descreve o funcionamento do
        circuito.
\end{listar}

Os modos das portas são \texttt{IN}, \texttt{OUT}, \texttt{BUFFER} e
\texttt{INOUT}. Uma porta \texttt{OUT} não pode ser lida dentro da
arquitetura, ao contrário de uma porta \texttt{BUFFER} \cite{aula3}. A
atribuição a um \texttt{SIGNAL} ou a uma porta usa o operador \texttt{<=}; os
comentários começam com \texttt{--} e vão até o fim da linha. A linguagem não
distingue maiúsculas de minúsculas, mas é recomendável escrever as palavras
reservadas em maiúsculas e os nomes de sinais em minúsculas \cite{aula3}.

O Código~\ref{lst:cap3ex1} reproduz o exemplo de código completo apresentado
na aula 3.

\break

\begin{lstlisting}[caption={Exemplo de ENTITY e ARCHITECTURE (aula 3).},label={lst:cap3ex1}]
ENTITY Cap3ex1 IS
    PORT (
        a     : IN BIT;
        b     : IN BIT;
        x,y,z : OUT BIT
    );
END Cap3ex1;
-----------------------------
ARCHITECTURE arch OF Cap3ex1 IS
BEGIN
    x <= a AND b;   -- saida x
    y <= a OR b;    -- saida y
    z <= a XOR b;   -- saida z
END arch;
\end{lstlisting}
\vspace{-0.3cm}
\fonte{ Codigo retirado do slide aula 6 de Garcia.}

\section{Representação em complemento de dois}

Em palavras de 8 bits em complemento de dois, o bit 7 é o bit de sinal \cite{harris2012}: valores de $-128$ a $127$ são representáveis, e o bit 7 igual a `1' indica número negativo. Por isso a flag de negativo (N) é simplesmente o bit 7 do resultado.

\section{Soma e subtração}

A soma bit a bit com propagação de carry é feita por um somador binário. A subtração pode ser reduzida a uma soma, pois
\begin{equation}
  a - b - c = a + \overline{b} + \overline{c},
  \label{eq:sub}
\end{equation}
onde $\overline{b}$ é o complemento bit a bit de $b$ e $\overline{c}$ é o inverso do carry de entrada. Isso permite que um único somador atenda às quatro operações aritméticas: basta inverter $b$ nas subtrações e escolher adequadamente o carry de entrada do somador (Tabela~\ref{tab:cin}).

\begin{table}[htb]
  \centering
  \caption{Carry de entrada do somador em cada operação aritmética.}
  \label{tab:cin}
  \begin{tabular}{llll}
    \hline
    Operação & Entrada B do somador & Carry de entrada do somador & Resultado \\ \hline
    ADDC & $b$ & $c_{in}$ & $a + b + c_{in}$ \\
    ADD  & $b$ & 0 & $a + b$ \\
    SUBC & $\overline{b}$ & $\overline{c_{in}}$ & $a - b - c_{in}$ \\
    SUB  & $\overline{b}$ & 1 & $a - b$ \\ \hline
  \end{tabular}
  \fonte{Autoria propria.}
\end{table}

\section{Flags de status}

\begin{description}
  \item[Carry/borrow (C)] Na soma, é o carry do bit 7. Na subtração, é o \textit{borrow} (empréstimo), igual ao carry do somador invertido. Nas rotações e deslocamentos guarda o bit que sai da palavra.
  \item[Half carry/borrow (H)] Análogo ao C, mas referente à passagem do bit 3 para o bit 4.
  \item[Zero (Z)] Vale `1' quando todos os bits do resultado são zero.
  \item[Negative (N)] Cópia do bit 7 do resultado.
  \item[Overflow (V)] Indica que o resultado em complemento de dois não cabe em 8 bits. Com $b_{op}$ sendo o operando efetivamente somado ($b$ ou $\overline{b}$), ocorre overflow quando os sinais de $a$ e $b_{op}$ são iguais e o sinal do resultado é diferente:
  \begin{equation}
    V = \overline{a_7 \oplus b_{op,7}} \cdot (a_7 \oplus r_7).
    \label{eq:ovf}
  \end{equation}
\end{description}

\secao{Operações aritméticas com std\_logic\_vector}\rotulo{sec:aritmetica}

O tipo \texttt{std\_logic\_vector} do pacote IEEE 1164 só permite operações
lógicas. Para realizar operações aritméticas há duas soluções \cite{aula6}:

\begin{listar}
  \item usar um dos pacotes \texttt{std\_logic\_unsigned} ou
        \texttt{std\_logic\_signed} (nunca os dois ao mesmo tempo);
  \item usar o pacote \texttt{numeric\_std}, que define os tipos
        \texttt{signed} e \texttt{unsigned} e funções de conversão como
        \texttt{to\_integer}, \texttt{to\_unsigned} e \texttt{to\_signed}. Esse
        pacote deve ser preferido ao pacote \texttt{arith}.
\end{listar}

Para detectar o \emph{carry} (``vai um'') no bit mais significativo, a operação
deve ser feita com 9 bits, e o bit de \emph{carry} é o bit mais significativo
do resultado \cite{aula6}. O Código~\ref{lst:alu3} reproduz o exemplo
\texttt{alu\_3} da aula 6, que implementa uma ALU de quatro operações com saída
de \emph{carry}, e serve de ponto de partida para o projeto.

\begin{lstlisting}[caption={Arquitetura do exemplo alu\_3 (aula 6).},label={lst:alu3}]
ARCHITECTURE arch OF alu_3 IS
    SIGNAL aux : STD_LOGIC_VECTOR(8 DOWNTO 0);
BEGIN
    WITH op SELECT
        aux <= '0' & (a AND b)        WHEN "00",
               '0' & NOT a            WHEN "01",
               ('0' & a) + ('0' & b)  WHEN "10",
               ('0' & a) - ('0' & b)  WHEN OTHERS;

    r <= aux(7 DOWNTO 0);
    c <= aux(8);
    z <= '1' WHEN aux(7 DOWNTO 0) = "00000000" ELSE '0';
END arch;
\end{lstlisting}
\vspace{-0.3cm}
\fonte{ Codigo retirado do slide aula 6 de Garcia.}
